\section{A chiral Pati--Salam block}
\label{sec:pstar}

The five projectors of \cref{sec:projectors} are one family among many. Enlarging the class to all commuting symmetric blade
involutions produces a projector on which colour and chirality coexist. The idea is to correlate spacetime chirality with an
\emph{internal} four-dimensional chirality, and to constrain the colour \emph{volume} rather than any colour direction.

\subsection{Construction}

At three constituents take the interior gamma matrices $\gamma^\mu=(e^{(1)}_1,e^{(1)}_4,e^{(1)}_5,e^{(1)}_6)$ and split the remaining
fourteen directions into three groups:
\begin{center}
\begin{tabular}{@{}llcc@{}}
\toprule
group & directions & signs & generators $E_A$\\
\midrule
weak $W$ & $e^{(2)}_1,\ e^{(3)}_1,\ e^{(3)}_2,\ e^{(3)}_3$ & $++++$ & $7,13,14,15$\\
colour $C$ & $e^{(2)}_4,e^{(2)}_5,e^{(2)}_6,\ e^{(3)}_4,e^{(3)}_5,e^{(3)}_6$ & $------$ & $10,11,12,16,17,18$\\
auxiliary $A$ & $e^{(1)}_2,e^{(1)}_3,\ e^{(2)}_2,e^{(2)}_3$ & $++++$ & $2,3,8,9$\\
\bottomrule
\end{tabular}
\end{center}
With $J=\wI$ the ordered product of the fourteen internal generators and $\Gamma=J\gamma^0\gamma^1\gamma^2\gamma^3$, define the three
volumes
\[
\chi_W=e^{(2)}_1e^{(3)}_1e^{(3)}_2e^{(3)}_3,\qquad
\chi_C=e^{(2)}_4e^{(2)}_5e^{(2)}_6e^{(3)}_4e^{(3)}_5e^{(3)}_6,\qquad
q_A=e^{(1)}_2e^{(1)}_3 .
\]

\begin{figure}[t]
\centering
\begin{tikzpicture}[scale=0.95]
\foreach \k in {1,2,3}{\node at (1.3*\k,6.55) {\small $k=\k$};}
\node[anchor=south] at (2.6,6.85) {\small constituent};
\foreach \a in {1,...,6}{\node at (0.35,6.5-\a) {\small $e_{\a}$};}
% colour code: interior = gray, W = gcol, C = ocol, A = tcol
\def\pscell#1#2#3{\fill[#3] (1.3*#1-0.5,6.0-#2) rectangle (1.3*#1+0.5,6.0-#2+0.9);\draw[gray] (1.3*#1-0.5,6.0-#2) rectangle (1.3*#1+0.5,6.0-#2+0.9);}
\pscell{1}{1}{lgray}\pscell{1}{4}{lgray}\pscell{1}{5}{lgray}\pscell{1}{6}{lgray}
\pscell{1}{2}{tcol!35}\pscell{1}{3}{tcol!35}
\pscell{2}{1}{gcol!40}\pscell{2}{2}{tcol!35}\pscell{2}{3}{tcol!35}
\pscell{2}{4}{ocol!45}\pscell{2}{5}{ocol!45}\pscell{2}{6}{ocol!45}
\pscell{3}{1}{gcol!40}\pscell{3}{2}{gcol!40}\pscell{3}{3}{gcol!40}
\pscell{3}{4}{ocol!45}\pscell{3}{5}{ocol!45}\pscell{3}{6}{ocol!45}
\foreach \k in {1,2,3}{\foreach \a in {1,...,6}{\pgfmathtruncatemacro{\A}{6*(\k-1)+\a}\node at (1.3*\k,6.45-\a) {\scriptsize $E_{\A}$};}}
\draw[thick] (0.75,3.05) -- (4.45,3.05);
\node[anchor=west,align=left] at (5.0,5.6) {\small \textcolor{mgray}{\rule{8pt}{8pt}} interior $\gamma^\mu$: Lorentz};
\node[anchor=west,align=left] at (5.0,4.95) {\small \textcolor{gcol!70}{\rule{8pt}{8pt}} weak $W$: volume $\chi_W$, $\su2_L\oplus\su2_R$};
\node[anchor=west,align=left] at (5.0,4.3) {\small \textcolor{ocol!80}{\rule{8pt}{8pt}} colour $C$: volume $\chi_C$, $\so6\cong\su4$};
\node[anchor=west,align=left] at (5.0,3.65) {\small \textcolor{tcol!70}{\rule{8pt}{8pt}} auxiliary $A$: $q_A=e^{(1)}_2e^{(1)}_3$};
\node[anchor=west,align=left] at (5.0,2.6) {\small $X=\Gamma\chi_W,\quad Z=-J\chi_C,\quad U=-Jq_A$};
\node[anchor=west,align=left] at (5.0,1.9) {\small $P_*=\tfrac18(\id+X)(\id+Z)(\id+U)$, rank $64$};
\node[anchor=west,align=left] at (5.0,1.2) {\small temporal rows above the line, spatial below};
\end{tikzpicture}
\caption{The eighteen generators of $\Cl(9,9)$ arranged by constituent, and the grouping that defines the Pati--Salam block. The
interior four-plane carries the Lorentz action; the weak, colour and auxiliary groups enter $P_*$ only through their volumes, which
is why every rotation inside each group survives the projection (\cref{thm:PS,thm:PSalg}).}
\label{fig:pstar}
\end{figure}

\begin{theorem}[A chiral projector carrying colour]
\label{thm:PS}
$X=\Gamma\chi_W$, $Z=-J\chi_C$ and $U=-Jq_A$ are commuting symmetric traceless involutions, and all seven nonidentity products $X$, $Z$,
$U$, $XZ$, $XU$, $ZU$, $XZU$ are traceless. Hence
\begin{equation}
\label{eq:Pstar}
P_*=\tfrac18(\id+\Gamma\chi_W)(\id-J\chi_C)(\id-Jq_A)
\end{equation}
is an orthogonal projector with $\Tr P_*=512/8=64$. It satisfies $[P_*,\Sigma^{\mu\nu}]=[P_*,J]=0$ for the six Lorentz generators
$\Sigma^{\mu\nu}=\tfrac12\gamma^\mu\gamma^\nu$, and $P_*\gamma^\mu P_*=0$, so it is axial; its two Weyl blocks have real rank $32$ each, and
$\chi_WP_*=\Gamma P_*$.
\end{theorem}

Rank is not automatic from the three factors being individually traceless; it is the products that decide, and here all of them are
traceless. The statement $\chi_WP_*=\Gamma P_*$ is the correlation of spacetime chirality with the internal chirality of the weak
directions.

\subsection{The effective algebra}

\begin{theorem}[The effective algebra and its Pati--Salam part]
\label{thm:PSalg}
All fifteen colour bivectors $e_ae_b$, $a,b\in C$, commute with $X$, $Z$ and $U$, so the whole compact $\so6_C\cong\su4_C$ survives
and acts faithfully on the block; the colour volume $\chi_C$ is invariant under those rotations, so constraining it costs no colour
generator. With $w_1,\dots,w_4$ the weak directions, $\tau_i=-\tfrac12w_jw_k$ for cyclic $(i,j,k)$, $t^{L,R}_i=\tfrac12(\id\mp\chi_W)\tau_i$
and $T^{L,R}_i=P_*t^{L,R}_i$,
\[
[T^L_i,T^L_j]=\epsilon_{ijk}T^L_k,\qquad T^L_iP_R=0,\qquad -\sum_i(T^L_i)^2=\tfrac34P_L,
\]
and the mirror identities hold on the right. The $t^{L,R}_i$ are genuine bivectors, for example $t^L_3=-\tfrac14(w_1w_2+w_3w_4)$.
The full effective compact bivector algebra is
\[
\keff(P_*)=\su4\oplus\su2_L\oplus\su2_R\oplus\uu1^2,
\]
of dimension $23$, entirely compact, with a derived algebra of dimension $21$ and a two-dimensional centre. All $21$ simple generators
preserve the kinetic kernels $K^\mu=P_*\gamma^0\gamma^\mu P_*$, and $K^0=P_*$.
\end{theorem}

\begin{proof}
Each colour bivector has even overlap with each of the three involution masks, so all fifteen commute with $X,Z,U$ and descend
faithfully. For the weak part, $\chi_W^2=\id$ and $[\chi_W,\tau_i]=0$ make $\tfrac12(\id\mp\chi_W)$ idempotents commuting with the
$\tau_i$; the $\tau_i$ close on $\su2$ with $\tau_i^2=-\tfrac14\id$, so $-\sum_i(T^L_i)^2=\tfrac34P_L$, and $\chi_WP_*=\Gamma P_*$ places the
two ideals on opposite chiralities. The dimension $23$ is the full solution of $[B,P_*]=0$ over the $91$-dimensional space of internal
bivectors: the nullspace is $23$-dimensional, entirely compact, has trivial annihilator on the block, and has a derived algebra of
dimension $21$.
\end{proof}

\begin{corollary}[One Pati--Salam family]
\label{cor:PSfamily}
The colour volume is not an involution: $\chi_C$ is a grade-six blade on six negative directions, so $\chi_C^2=-\id$. The
involution is
\[
\Gamma_C:=J\chi_C,\qquad \Gamma_C^2=\id,\qquad \Gamma_CP_*=-P_*,\quad\text{equivalently}\quad \chi_CP_*=JP_*.
\]
On the module with $J=+i$ each Weyl block then carries exactly four $\su4$ weights, all of colour chirality $\Gamma_C=-1$, each with
multiplicity $4=2\times2$, a weak doublet times a Weyl spinor, and the two blocks carry the same colour weights. The content is
\[
(\mathbf4,\mathbf2,\mathbf1)_L\oplus(\mathbf4,\mathbf1,\mathbf2)_R ,
\]
one Pati--Salam family including a right-handed neutrino \cite{PatiSalam}. Which of the two conjugate four-dimensional representations
is called $\mathbf4$ is a convention.
\end{corollary}

\subsection{Charges}

The compact algebra is realised by skew operators; charges are the corresponding symmetric operators $\hat t=-Jt$. On the block the
centre of $\keff(P_*)$ is spanned by the skew operators $JP_*$ and $J\Gamma P_*$, whose charges are $P_*$ and $\Gamma P_*$. Concretely
$q_A|_{V_*}=J$ and $E_8E_9|_{V_*}=J\Gamma$ on the image $V_*=\operatorname{im}P_*$.

\begin{proposition}[Colour intertwiner, $\BmL$, and the charge table]
\label{prop:PSledger}
$R=P_*\gamma^0e^{(2)}_1e^{(2)}_2$ satisfies $R^2=-P_*$, exchanges the two Weyl blocks and commutes with every colour generator, so the
colour representations of the two chiralities are equivalent and the content is $(\mathbf4,\mathbf2,\mathbf1)_L\oplus(\mathbf4,\mathbf1,
\mathbf2)_R$ rather than a $\mathbf4/\bar{\mathbf4}$ pairing. With
\[
H=-\bigl(e^{(2)}_4e^{(3)}_4+e^{(2)}_5e^{(3)}_5+e^{(2)}_6e^{(3)}_6\bigr),\qquad Q_c=-JH,
\]
the derived centraliser of $H$ inside colour $\so6$ is $\su3$, and $(Q_c-\id)(Q_c+3\id)P_*=0$: $Q_c$ has eigenvalues $+1$ and $-3$ on
the block, with real multiplicities $48$ and $16$, the triplet and the singlet. Setting $\BmL=Q_c/3$ and using the weak charges
$-JT^3_{L,R}$,
\[
Y=-J\,T^3_R+\tfrac12(\BmL),\qquad Q=-J\,T^3_L+Y ,
\]
gives one family:
\begin{center}
\begin{tabular}{@{}lccr@{}}
\toprule
field & $SU(3)\times SU(2)_L$ & $Y$ & internal Weyl components\\
\midrule
$Q_L$ & $(\mathbf3,\mathbf2)$ & $\tfrac16$ & $6$\\
$L_L$ & $(\mathbf1,\mathbf2)$ & $-\tfrac12$ & $2$\\
$u_R$ & $(\mathbf3,\mathbf1)$ & $\tfrac23$ & $3$\\
$d_R$ & $(\mathbf3,\mathbf1)$ & $-\tfrac13$ & $3$\\
$e_R$ & $(\mathbf1,\mathbf1)$ & $-1$ & $1$\\
$\nu_R$ & $(\mathbf1,\mathbf1)$ & $0$ & $1$\\
\bottomrule
\end{tabular}
\end{center}
Four real spinor coordinates make one internal Weyl component, so rank $64$ is sixteen internal Weyl components, eight left and eight
right, and not sixty-four particles.
\end{proposition}

The $\su3$ embedding, the right-handed Cartan direction and the hypercharge combination are chosen; what is established is their
mutual compatibility. In the complex structure $J$ every one of the sixteen Weyl fermions has charge one: $J$ is fermion number,
$F=3B+L$, and the baryon number appears as $4B=F+(\BmL)$.

\begin{proposition}[The two central factors cannot be gauged]
\label{prop:PSanom}
For a central charge $a+b\Gamma$, corresponding to the skew element $aJ+bJ\Gamma$, the anomaly coefficients on the block, per Weyl
fermion with every field written left-handed, are
\[
\begin{gathered}
SU(2)_L^2:\ 2(a-b),\qquad SU(2)_R^2:\ -2(a+b),\qquad SU(4)^2:\ -2b,\\
\text{cubic}:\ -16(3a^2b+b^3),\qquad \text{gravitational}:\ -16b ,
\end{gathered}
\]
and the mixed anomalies with two central and one Pati--Salam generator vanish. All vanish only for $a=b=0$. The Pati--Salam algebra itself
is anomaly free on the block: the $SU(4)^3$ anomaly cancels between two fundamentals and two antifundamentals in the all-left convention,
and each weak $SU(2)$ sees four doublets, so the global anomaly is absent \cite{Witten82}. Hence, with the block as the matter content,
$\su4\oplus\su2_L\oplus\su2_R$ is the largest anomaly-free gauge algebra containing Pati--Salam, and the two central directions remain
global symmetries.
\end{proposition}

The obstruction is stated for this matter content; additional matter or another cancellation mechanism is not excluded. The two central
symmetries are fermion number $J$, anomalous only through $SU(2)_L$ and $SU(2)_R$ and hence exact in perturbation theory and violated by
their instantons \cite{tHooft76}, and the chiral charge $J\Gamma$.

\begin{proposition}[Two rank-64 projectors, two orbits]
\label{prop:orbits}
Under the compact group of the connection sector, whose Lie algebra is the maximal compact $\so8\oplus\so6$ of dimension $43$, the orbit
through $\Pfam$ has dimension $43-9=34$ and the orbit through $P_*$ has dimension $43-23=20$. Equal rank does not put two projectors on
one orbit.
\end{proposition}

\begin{remark}[What the construction settles]
\label{rem:PSsettles}
It settles an existence question: colour on a chiral block is not obstructed, and no counting argument based on the five-projector family
can be a theorem. It does not select $P_*$; it does not derive the hypercharge embedding, which is compatible with anomaly cancellation but
chosen; it gives one family and not three; and by \cref{cor:nopredict} exhibiting a known representation content is a consistency result,
not evidence. One further fact will matter in \cref{sec:matter}: $\gamma^0$ commutes with all $91$ internal bivectors and with $J$ and the
rotations, and anticommutes with $\Gamma$ and the boosts. It maps $V_*$ onto an orthogonal rank-$64$ subspace on which the same internal
algebra acts with $\Gamma$ reversed. The chirality of the block is put in by the choice of $P_*$.
\end{remark}
